% =================================================
% C 册补充专题：男性作为性侵受害者（呼应 A 册《性侵犯与强奸》）
% 现代循证撰写
% =================================================

\chapter{男性作为性侵受害者（专科补充）}

\section{先把"男人也该扛住"放一边}
性侵不只是女性会遇到的事。\textbf{男性、男性之间的、以及异性恋男性被女性施加的}性侵都存在，
只是长期被"男人要主动、要硬、要能行"的刻板印象遮住，导致男性受害者\textbf{更少求助、更晚被发现}。
对男性来说，最大的障碍往往不是疼痛，而是\textbf{羞耻与被质疑}（"一个大男人怎么会""你当时怎么不反抗"）。
请把这两句先收下：\textbf{你是受害者，与你的性别、体力、是否"想发生"都无关；求助不是软弱。}

\section{男性特有的医学要点}
\begin{itemize}
  \item \textbf{生殖道 / 泌尿道损伤}：尿道或阴茎撕裂、血尿、射精痛，可能伤及尿道海绵体；出现"排尿困难、血尿、阴囊肿痛"要尽快到\textbf{急诊 / 泌尿外科}。
  \item \textbf{STI 暴露}：与 A 册一致——HIV（72 小时内 PEP 最有效）、梅毒、淋病、HPV、肝 B 等，需检测与随访。
  \item \textbf{意外致他人妊娠}：若为阴道 / 肛交且无保护，\textbf{伴侣}可能怀孕——这是男性受害者常被忽略、却又必须面对的一环。
  \item \textbf{会阴疼痛、盆底肌紧张}：部分男性事后出现会阴痛、勃起或射精不适，多与盆底肌痉挛 + 心理因素相关，可治疗。
\end{itemize}

\section{就医时怎么说（给"开不了口"的男性）}
\begin{itemize}
  \item 直接说：\textbf{"我可能遭遇了性侵，想检查一下伤情和感染风险，请帮我做相关检查。"}——一句话就够，不必展开细节。
  \item 可要求：体格检查 + 分泌物/尿液取样 + 血检（HIV/梅毒/乙肝）+ 必要时 STI 随访 + 伴侣妊娠评估的指引。
  \item 若不愿走法律程序，可只体检取证、保留证据，之后再决定是否报案；身体安全永远第一。
  \item 找\textbf{男性医生}或\textbf{同性别医护人员}陪同检查，能明显降低不适感。
\end{itemize}

\section{心理与社会支持}
\begin{itemize}
  \item 男性受害者同样可能出现 PTSD、闪回、对亲密关系退缩、性欲骤变（可能升高也可能降低，都正常）。
  \item 支持渠道：心理科 / 创伤聚焦咨询（TF-CBT、EMDR）、男性权益或心理援助热线；必要时伴侣共同咨询。
  \item 一句给男性的话：\textbf{"扛住"不是勇敢，"说出口"才是；你不必为别人的暴力负责。}
\end{itemize}

\section{一页要点}
男性也是性侵受害者，路径与 A 册一致：安全→就医（泌尿外科/急诊）→STI/损伤/伴侣妊娠评估→心理支持。
男性专科要点：生殖道与泌尿道损伤、会阴/盆底疼痛、致伴侣妊娠风险。
开不了口时，一句话"我可能遭了性侵，想检查"就够；找同性别医护陪同、按自己节奏来。


